\documentclass{article}
\usepackage[a4paper, total={6.5in, 10in}]{geometry}
\setlength{\parindent}{0pt}
\usepackage{tcolorbox}
\tcbset{colback=white!94!black, colframe=black!60!white, coltitle = white, boxrule=0.8pt, arc=2mm}
\newtcolorbox{boxs}[2][]{title= #2,#1}
\usepackage{datetime2}
\usepackage[british]{babel}
\usepackage{amsfonts}
\usepackage{amsmath}

\title{\textbf{Probablity Lecture 1}}
\author{Robert O'Connell}
\date{\today}
\begin{document}
\maketitle

\subsection*{Sample Space}
A set of all possible outcomes(\(\Omega\))

Set must be
\begin{itemize}
    \item Mutually exclusive
    \item Collectively exhaustive
    \item At the "right" granularity
\end{itemize}

\subsection*{Probablity Axioms}

\textbf{Event}: A subset of the sample space

Probablities are assigned to events
\\

The Axioms are
\begin{enumerate}
    \item Nonnegativity \(P(A) \geq 0\)
    \item Normalization \(P(\Omega) = 1\)
    \item (Finite) Additivity, If \(A \cap B = \emptyset\) then \(P(A \cup B ) = P(A) + P(B)\)
\end{enumerate}

\subsubsection*{Simple Consequences of the Axioms}
We know,
\[
A \cup A^c = \Omega
\]
and,
\[
A \cap A^c = \emptyset
\]

Thus from our axioms we get,
\[
1 = P(\Omega) = P(A \cup A^c) = P(A) + P(A^c)
\]
Also,
\[
P(A) = 1 - P(A^c) \leq 1
\]

Finally,
\[
1 = P(\Omega) + P(\Omega^c)
\]
\[
1 = 1 + P(\emptyset)
\]
\[
P(\emptyset) = 0
\]

\subsubsection*{Strengthining Axiom 3}
\[
P(A \cup B \cup C) = P((A \cup B) \cup C) = P(A \cup B) + P(C) = P(A) + P(B) + P(C)
\]

Thus we can see that,

If \(A_1, \dots , A_k\) are disjoint events \(\Rightarrow P(A_1 \cup \dots \cup A_k) = \sum_{i=1}^{k} P(A_i) \)

\subsubsection*{Further Consequences}
Let \(A \subset B\),
\[
B = A \cup (B \cap A^c)
\]
\[
P(B) = P(A) + P(B \cap A^c) \geq P(A)
\]
\\

Let A, B be two sets such that \(P(A \cap B) \neq \emptyset\)

We want \(P(A \cup B) = P(A) + P(B) + P(A \cap B)\)

Let,
\[
a = P(A \cap B^c), \quad b = P(A \cap B), \quad c = P(B \cap A^c)
\]
Then,
\[
P(A \cup B) = a + b + c
\]
\[
P(A) + P(B) - P(A \cap B) = (a+b) + (b+c) - b = a + b + c
\]
\\

Since \(P(A \cap B) \geq 0\), \(P(A \cup B) \leq P(A) + P(B)\)   (union bound)
\\

Finally,
\[
A \cup B \cup C = A \cup (B \cap A^c) \cup (C \cap A^c \cap B^c)
\]
These events are disjoint, thus
\[
P(A \cup B \cup C) = P(A) + P(B \cap A^c) + P(C \cap A^c \cap B^c)
\]

\subsection*{Discrete Uniform Law}
Assume \(\Omega\) contains \(n\) equally likely elements, assume the event A
consists of \(k\) elements.

Then
\[
P(A) = k \cdot \frac{1}{n}
\]
\\

If \(\Omega \) is not discrete, Probablity = Area

\subsection*{Probablity Calculation Steps}
\begin{itemize}
    \item Specify the sample space
    \item Specify a probabilty law
    \item Identify an event of interest
    \item Calculate...
\end{itemize}

\subsection*{Countable Additivity Axiom}
If \(A_1,A_2,A_3 \dots\) is an infinte sequence of disjoint events then

\(P(A_1 \cup A_2 \cup A_3 \cup \dots) = P(A_1) + P(A_2) + P(A_3) + \dots\)

\subsubsection*{Note}
This only holds for "countable" sequences of events (shouldn't be a problem in this module)
\\

Using this Axiom on the unit square we see that it doesn't hold
\[
1 = P(\Omega) = P(\cup \{(x,y)\}) = \sum P(\{(x,y)\}) = \sum 0 = 0
\]
\end{document}